\addchap{Abbreviations and conventions}\label{ch:Abb}

Phonetic transcription and characters generally follow IPA with the following exceptions.
For high front semi-vowel, we use ‹y› rather than ‹j›
as is common in eastern Indonesian and Austronesian studies.
Similarly, the transcription of Austronesian reconstructions
follows conventions that are widely used in that field, as summarised in §1.7.

\section*{Symbols, punctuation, and styles}

{\wg
\begin{xltabular}{\textwidth}[l]{l>{\raggedright\arraybackslash}p{5cm}>{\raggedright\arraybackslash}X}
\lsptoprule
	&	{use}	&	{comment or examples}	\\ \midrule \hhline{~}
\endhead
§	&	cross reference, section, see section 	&	(§1.4.1, §11.1.3)	\\
\tcg{(text)}	&	non-cognate forms in data tables	&	*duRi, duri, luri, ruri, \tcg{(isinan)}, luli	\\
{[text]}	&	phonetic data	&	{[ˈhaʤu]}, {[ˈʧia]}	\\
/text/	&	phonemic data, underlying morphemes	&	/en-suba-n/, /pa-dətu/	\\
‹text›	&	orthographic, typographic	&	‹soë›, ‹matèj›, ‹ràmma›	\\ %‹›
{\<}  {\>}	&	infix, infixation	&	*b\<in\>ahi, \it{w\<i\>ara}	\\ %
>	&	becomes, derives 	&	PMP *qulu > \it{olo-n} ‘head’; *p > \it{h}	\\
<	&	comes from, derives from	&	\it{uki} < PMP *punti ‘banana’	\\
<	&	less than	&	<4,000 years	\\
*	&	reconstructed proto phoneme; 	&	*b, *q, *h, *ə, *-aw;	\\
*	&	reconstructed etymon (PMP)	&	*babuy ‘pig’, *pija ‘how much?’	\\
**	&	reconstruction at node below PMP (proto phoneme or etymon)	&	**dama ‘eye’, **maya ‘tongue’	\\
*[   ]	&	unattested, not found, ungrammatical	&	*{[ubu]}	\\
=	&	historical retention 	&	*p = \it{p}, *d = \it{d}, *s = \it{s}	\\
/	&	merger	&	*z/*y > **y; *a/*-ay > \it{a}	\\
{,} \phantom{ }	&	split	&	*j > \it{l}, \it{r}; *p > \it{f, h}	\\ \hhline{~}
\phantom{ }	&	\phantom{ }	&	(most common reflex listed first)	\\
{\0}	&	loss of historical phoneme or affix	&	*q > {\0}; *j > \it{l}, {\0}; *p > \it{f}, \it{h}, {\0}	\\
{\0}\tsc{+fr}	&	loss with fronting of adjacent segment	&	*y > {\0}\tsc{+fr}	\\
{[LS]}	&	lexically specific (split, merger)	&	*j > \it{l}, {\0} [LS]; *p > \it{f, h, {\0}} {[LS]}	\\
{\#}	&	pseudo-reconstruction (found in unrelated language families), or ad-hoc reconstruction	&	{\#}muku ‘banana’; {\#}kisor(o) ‘cuscus’	\\
{\#}	&	word boundary; initial or final	&	\it{{\#}b} = initial \it{b}; \it{p{\#}} = final \it{p} at end of word	\\
C	&	consonant	& \phantom{ }\\
CC	&	consonant cluster	& \phantom{ }	\\
V	&	vowel	& \phantom{ }	\\
VV	&	vowel sequence	&	can be same (like vowels; \it{aa, oo}), or different (unlike; \it{ao, oe})	\\
NC	&	nasal-consonant cluster	&	 \phantom{ }\\
σ	&	syllable	&	\phantom{ }	\\
ˈ	&	stress	&	ˈCVCVC; ˈσσ; {[ˈʔəm:u]} ‘house’	\\
{\Q}	&	orthographic glottal or implosive	&	Nga{\Q}o, Palu{\Q}e, ‹b{\Q}› = {[ɓ]}	\\
{\tl}	&	varies with	&	\it{leya {\tl} leiya {\tl} lea}	\\
{\tl}	&	reduplication boundary	&	\it{gos{\tl}gosa; emsi{\tl}emsikan}	\\
→	&	synchronic realisation	&	/b/ → {[b]} {\tl} {[β]}	\\
{ ̆ }&	extra short vowel	&	ă ĕ ĭ ŏ ŭ ə̆	\\
ê	&	\tsc{Aru} mid-high front vowel contrasting with mid-low [ɛ]		& \phantom{ }\\ 
ô	&	\tsc{Aru} mid-high front vowel contrasting with mid-low [ɔ]		& \phantom{ }\\
{[{-}{-}{-}]}	&	no ISO code	&	\ili{Sanggar} [{-}{-}{-}]	\\
{[xxx]}	&	primary ISO variety	&	or best documented variety	\\
(xxx)	&	secondary ISO variety	&	may need separate code; regular differences in sound correspondences and vocabulary	\\
\tsc{+fr}	&	\tsc{+front}	&	raising or fronting of V, or palatalisation of C	\\
\tsc{tr}	&	traces found (broader than \tsc{+fr})	&	can include sporadic *R > \it{ʔ}, etc.	\\
*a(C)i	&	secondary diphthong	&	from loss of historical C	\\
*gloss	&	gloss of proto-form	&	*duRi ‘thorn’	\\
local gl.&	gloss of reflex in local languages	&	*duRi ‘bone’	\\
\tsc{small}	&	subgroup or node with other & \tsc{\ili{Bima}-Lembata}, \tsc{Timor-Babar}, \\ \hhline{~}
\tsc{caps}	& nodes or languages below it & \tsc{Rote-Meto}, \tsc{SW Maluku} \\ \hhline{~}
\phantom{ } & \phantom{ } &	\tsc{Tanimbar-Bomberai}	\\
\tbr{red}	&	\mc{2}{l}{segments under discussion or in focus}		\\
\tbb{blue}	&	\mc{2}{l}{second type of segment in focus}		\\
\lspbottomrule
\end{xltabular}
}

\section*{Glosses}

{\wg
\begin{xltabular}{\textwidth}[l]{l>{\raggedright\arraybackslash}p{5cm}>{\raggedright\arraybackslash}X}
\lsptoprule
	gloss	&	expansion	\\ \midrule \hhline{~}
\endhead
1	&	first person	\\	
2	&	second person	\\	
3	&	third person	\\	
A	&	Actor, macrorole	\\	
\tsc{art}	&	article	\\	
\tsc{caus}	&	causative	\\	
\tsc{dist}	&	distal in time, space, or reference; that, there, then	\\	
\tsc{dl}	&	dual	\\	
\tsc{excl}	&	exclusive	\\	
\tsc{gen}	&	genitive	\\	
\tsc{hum}	&	human 	\\	
\tsc{incl}	&	inclusive	\\	
\tsc{ipfv}	&	imperfective	\\	
\tsc{irr}	&	irrealis: intended, purpose, future, about to, going to	\\	
\tsc{loc}	&	locative: in, at, on	\\	
\tsc{neg}	&	negation, negative, no, not	\\	
\tsc{neg1}	&	first part of bipartite negation	\\	
\tsc{neg2}	&	second part of bipartite negation	\\	
\tsc{nhum}	&	non-human 	\\	
\tsc{nmlz}	&	nominalizer	\\	
\tsc{pfv}	&	perfective ‘already’	\\	
\tsc{pi}	&	plural inclusive	\\	
\tsc{pl}	&	plural 	\\	
\tsc{poss}	&	possessive	\\	
\tsc{prox}	&	proximal in time, space or reference; this, here, now	\\	
\tsc{px}	&	plural exclusive	\\	
\tsc{recip}	&	reciprocal	\\	
S\sub{A}	&	single argument of active intransitive clause (subject as Actor)	\\	
\tsc{sg}	&	singular 	\\	
\tsc{stat}	&	stative (BE relationship)	\\	
S\sub{U}	&	single argument of non-active intransitive clause (subject as Undergoer)	\\	
U	&	Undergoer, macrorole	\\	\lspbottomrule
\end{xltabular}}

\section*{Subgroups}\label{sec:Subgroups}

{\wg
\begin{tabular}{ll}
\lsptoprule				
{abbreviation}	&	{expansion}	\\	\midrule
\tsc{An}	&	Austronesian	\\	
AP	&	\ili{Alor}-\ili{Pantar} (node under TAP; Papuan languages)	\\	
AS	&	\ili{Ambon}-Seram	\\	
BL	&	\ili{Bima}-Lembata	\\	
\tsc{Cel}	&	Celebic	\\	
CEMP	&	\ili{Central}-Eastern Malayo-Polynesian \it{(disputed node)}	\\	
CM	&	“\ili{Central} Maluku” \it{(geographical grouping)}	\\	
CMP	&	\ili{Central} Malayo-Polynesian \it{(disputed node)}	\\	
CT	&	\ili{Central} Timor	\\	
EMP	&	Eastern Malayo-Polynesian	\\
FL	&	Flores-Lembata (node of \ili{Bima}-Lembata) \\
GWB	&	Greater West Bomberai (Papuan, possibly linked to TNG)	\\	
MP	&	Malayo-Polynesian	\\	
non-\tsc{An}	&	non-Austronesian, unrelated Papuan language families	\\	
\tsc{Oc}	&	Oceanic	\\	
SB	&	Sula-\ili{Buru}	\\	
SHWNG	&	South Halmahera-West New Guinea	\\	
\tsc{SSul}	&	South Sulawesi	\\	
STB	&	Seram-Tanimbar-Bomberai	\\
\tsc{Tal}	&	\ili{Taliabu}	\\	
TAP	&	Timor-\ili{Alor}-\ili{Pantar} (Papuan language family)	\\
TB	&	Timor-\ili{Babar}	\\	
TNG	&	Trans-New Guinea (Papuan language family)	\\	
TNS	&	\ili{Teun}-\ili{Nila}-\ili{Serua}	\\	
WMP	&	Western Malayo-Polynesian \it{(disputed node)}	\\	\lspbottomrule
\end{tabular}
}

\section*{Languages}\label{sec:Languages}
{\wg
\begin{tabular}{ll}
\lsptoprule				
{abbreviation}	&	{expansion}	\\	\midrule
Anakal.	&	\ili{Anakalang}	\\
\ili{Bobot} (At.)	&	\ili{Atiahu} \ili{Bobot}	\\
\ili{Bobot} (We.)	&	Werinama \ili{Bobot}	\\
\ili{Dela}-O.	&	\ili{Dela}-\ili{Oenale}	\\
Erokw.	&	\ili{Erokwanas}	\\
\ili{Geser}-G.	&	\ili{Geser}-\ili{Gorom}	\\
Kamb.	&	\ili{Kambera}	\\
\ili{Keo}, Ud.	&	Udiworowatu \ili{Keo}	\\
Kusa-M.	&	\ili{Kusa-Manea}	\\
Lewo El.	&	\ili{Lewo Eleng}	\\
Liana-S.	&	\ili{Liana-Seti}	\\
Mamb.	&	\ili{Mamboru}	\\
\ili{Nga'o}, Wat.	&	Watumite \ili{Nga'o}	\\
\ili{Nila} (W)	&	\ili{Nila} from Wotay village	\\
\ili{Nila} (K)	&	\ili{Nila} from Kokroman village	\\
Uruang.	&	\ili{Uruangnirin}	\\
Watub.	&	\ili{Watubela}	\\
\ili{Wetan} (T)	&	\ili{Wetan} from \ili{Teun} Island	\\
\ili{Wetan} (W)	&	\ili{Wetan} from \ili{Wetan} island	\\
Wey.	&	\ili{Weyewa}	\\
\lspbottomrule
\end{tabular}
}



\newpage
\section*{Proto languages and Historical reconstructions}
Unless otherwise specified, all reconstructions are from \citet{BlustTrussel2020}.
PMP *wahiyəR ‘water’ is from \citet[1027]{Wolff2010}, who reconstructs PAN *wasiyeɣ.

{\wg
\begin{xltabular}{\textwidth}[l]{l>{\raggedright\arraybackslash}X}
\lsptoprule				
{abbreviation}	&	{expansion}	\\	\midrule \hhline{~}
\endhead
PAMS	&	Proto-\ili{Ambon}, \citet{Stresemann1927} = PCM	\\	
\tsc{PAn}	&	Proto-Austronesian	\\	
\tsc{PAnD}	&	Proto-Austronesian = PMP, \citet{Dempwolff1938} \\	
\tsc{PAn-F}	&	Proto-Austronesian (Formosan)	\\	
\tsc{PAnS}	&	Proto-Austronesian = PMP, \citet{Stresemann1927}	\\	
PAS	&	Proto-\ili{Ambon}-Seram	\\	
\it{PCEMP}	&	Proto-\ili{Central}-Eastern Malayo-Polynesian \it{(disputed node)}	\\	
P(CE)MP	&	PMP or PCEMP	\\	
PCM	&	Proto-\ili{Central} Maluku	\\	
\it{PCMP}	&	Proto-\ili{Central} Malayo-Polynesian \it{(disputed node)}	\\	
%PECM	&	Proto-East \ili{Central} Maluku	\\	
PEMP	&	Proto-Eastern Malayo-Polynesian	\\	
PGWB	&	Proto-Greater West Bomberai \citep{UsherSchapper2022}	\\	
PMP	&	Proto-Malayo-Polynesian	\\	
\tsc{POc}	&	Proto-Oceanic	\\	
\tsc{PPh}	&	Proto-Philippines	\\
PRM	&	Proto-Rote-\ili{Meto} \citep{Edwards2021}	\\	
%PSB	&	Proto-Sula-\ili{Buru}	\\	
%PTAP	&	Proto-Timor-\ili{Alor}-\ili{Pantar} (Papuan)	\\	
%PTNS	&	Proto-\ili{Teun}-\ili{Nila}-\ili{Serua}	\\	
\it{PWMP}	&	Proto-Western Malayo-Polynesian \it{(disputed node)}	\\	
P(W)MP	&	PMP or PWMP	\\	\lspbottomrule
\end{xltabular}}

\section*{Other}

{\wg
\begin{xltabular}{\textwidth}[l]{l>{\raggedright\arraybackslash}X}
\lsptoprule			
{abbreviation}	&	{expansion}	\\	\midrule
\endhead
a.k.a.	&	also known as	\\
\tsc{ad}	&	after Christ (in historical dating; Latin: \it{anno domini})	\\
\tsc{bc}	&	before Christ (in archaeological or historical dating)	\\
\tsc{bce}	&	before common era (in archaeological or historical dating)	\\
\tsc{bp}	&	before present (archaeological dating)	\\
C	&	child (C-in-law [son or daughter-in-law]; MBC = mother’s brother’s child etc.)	\\
C.	&	\ili{Central} (e.g. C. Timor)	\\
cal	&	calibrated (archaeological dating)	\\
CC	&	grandchild (child's child)	\\
cf.	&	compare; see also	\\
E.	&	\ili{east}	\\
e.g.	&	for example	\\
env.	&	environment	\\
eSib	&	elder sibling	\\
et al.	&	and others (in multi-authored bibliographical references)	\\
etc.	&	et cetera; and so forth (other examples not presented)	\\
f & and following page (e.g. 2008:75f = 2008:75–76) \\
ff & and following page (e.g. 2008:75ff) \\
(f.s.)	&	female speaking (e.g. *ñaRa ‘brother (f.s.)’)	\\
FZ	&	father's sister	\\
Glotto.	&	Glottocode \citep{Hammarstrom2023}	\\
how m.?	&	how much/many?	\\
i.e.	&	in other words	\\
IPA	&	International Phonetic Alphabet	\\
irr.	&	irregular (sound change)	\\
ISEA	&	Island South\ili{east} Asia	\\
ISO	&	International Standards Organisation; unique language identifier protocols 639-3 	\\
lit:	&	literally	\\
\it{(m.)}	&	metathesis	\\
(m.s.)	&	male speaking (e.g. *bətaw ‘sister (m.s.)’)	\\
MB	&	mother's brother	\\
\it{(met.)}	&	metathesis	\\
%Ml.	&	\ili{Malay} (abbreviation used in \citealt{vanEkris1864,vanEkris1865})	\\
Mly.	&	\ili{Malay}	\\
MSEA	&	Mainland South\ili{east} Asia	\\
(n.)	&	noun, nominal (e.g. sail (n.))	\\
N.	&	North (e.g. N. \ili{Wemale})	\\
pers. comm.	&	personal communication	\\
PP	&	grandparent (parent's parent)	\\
reg.	&	regular (sound change)	\\
SE.	&	South\ili{east} (e.g. SE Sulawesi)	\\
SOV	&	Subject-Object-Verb	\\
sp.	&	species (singular; flora and fauna)	\\
spp.	&	species (plural; flora and fauna)	\\
%Sul.	&	Sulawesi	\\
SVO	&	Subject-Verb-Object	\\
SW.	&	Southwest	\\
(v.)	&	verb (e.g.  sail (v.))	\\
vi.	&	intransitive verb	\\
VOC	&	Vereenigde Oost-Indische Compagnie (Dutch East Indies Trading Company)	\\
VSO	&	Verb-Subject-Object	\\
vt.	&	transitive verb	\\
W.	&	west	\\
WB	&	wife’s brother	\\
ySib	&	younger sibling	\\
ZH	&	sister’s husband	\\
\lspbottomrule
\end{xltabular}}

\addchap{Alphabetical list of abbreviations}

{\wg\st{3pt}
\begin{xltabular}{\textwidth}[l]{l>{\raggedright\arraybackslash}X}
\lsptoprule			
{abbreviation}	&	{expansion}	\\	\midrule
\endhead
A	&	Actor, macrorole	\\
\tsc{An} & Austronesian \\
a.k.a.	&	also known as	\\
\tsc{ad}	&	after Christ (in historical dating; Latin: anno domini)	\\
Anakal.	&	\ili{Anakalang}	\\
AP	&	\ili{Alor}-\ili{Pantar} (node under TAP; Papuan languages)	\\
\tsc{art}	&	article	\\
AS	&	\ili{Ambon}-Seram	\\
\tsc{bc}	&	before Christ (in archaeological or historical dating)	\\
\tsc{bce}	&	before common era (in archaeological or historical dating)	\\
BL	&	\ili{Bima}-Lembata	\\
\ili{Bobot} (At.)	&	\ili{Atiahu} \ili{Bobot}	\\
\ili{Bobot} (We.)	&	Werinama \ili{Bobot}	\\
\tsc{bp}	&	before present (archaeological dating)	\\
C	&	child (C-in-law [son or daughter-in-law]; MBC = mother’s brother’s child)	\\
C	&	consonant	\\
C.	&	\ili{Central} (e.g. C. Timor)	\\
cal	&	calibrated (archaeological dating)	\\
\tsc{caus}	&	causative	\\
CC	&	consonant cluster	\\
CC	&	grandchild (child's child)	\\
\tsc{Cel}	&	Celebic	\\
CEMP	&	\ili{Central}-Eastern Malayo-Polynesian \it{(disputed node)}	\\
cf.	&	compare; see also	\\
CM	&	“\ili{Central} Maluku” \it{(geographical grouping)}	\\
CMP	&	\ili{Central} Malayo-Polynesian \it{(disputed node)}	\\
CT	&	\tsc{\ili{Central} Timor}	\\
\ili{Dela}-O.	&	\ili{Dela}-\ili{Oenale}	\\
\tsc{dist}	&	distal in time, space, or reference; that, there, then	\\
\tsc{dl}	&	dual	\\
E.	&	\ili{east}	\\
e.g.	&	for example	\\
EMP	&	Eastern Malayo-Polynesian	\\
env.	&	environment	\\
Erokw.	&	\ili{Erokwanas}	\\
eSib	&	elder sibling	\\
et al.	&	and others (in multi-authored bibliographical references)	\\
etc.	&	etcetera; and so forth (other examples not presented)	\\
\tsc{excl}	&	exclusive	\\
FL	&	Flores Lembata (node of \ili{Bima}-Lembata) \\
(f.s.)	&	female speaking (e.g. *ñaRa ‘brother (f.s.)’)	\\
\tsc{fr}	&	\tsc{front} (e.g. V\tsc{+fr} = \it{i, e})	\\
FZ	&	father's sister	\\
\tsc{gen}	&	genitive	\\
\ili{Geser}-G.	&	\ili{Geser}-\ili{Gorom}	\\
gl.	&	gloss	\\
Glotto.	&	Glottocode \citep{Hammarstrom2023}	\\
GWB	&	Greater West Bomberai (Papuan, possibly linked to TNG)	\\
\tsc{h}	&	\tsc{high} (e.g. V\tsc{[+high]} = \it{i, u})	\\
how m.?	&	how much/many?	\\
\tsc{hum}	&	human 	\\
i.e.	&	in other words	\\
\tsc{incl}	&	inclusive	\\
IPA	&	International Phonetic Alphabet	\\
\tsc{ipfv}	&	imperfective	\\
\tsc{irr}	&	irrealis: intended, purpose, future, about to, going to	\\
ISEA	&	Island South\ili{east} Asia	\\
ISO	&	International Standards Organisation; unique language identifier protocols 639-3 	\\
Kamb.	&	\ili{Kambera}	\\
\ili{Keo}, Ud.	&	Udiworowatu \ili{Keo}	\\
Kusa-M.	&	\ili{Kusa-Manea}	\\
Lewo El.	&	\ili{Lewo Eleng}	\\
Liana-S.	&	\ili{Liana-Seti}	\\
lit:	&	literally	\\
\tsc{loc}	&	locative: in, at, on	\\
{[LS]}	&	lexically specific (split, merger)	\\
\it{(m.)}	&	metathesis	\\
Mamb.	&	\ili{Mamboru}	\\
MB	&	mother's brother	\\
\it{(met.)}	&	metathesis	\\	
%Ml.	&	\ili{Malay} (abbreviation used in \citealt{vanEkris1864,vanEkris1865})	\\	
Mly.	&	\ili{Malay} \\	
MP	&	Malayo-Polynesian	\\
(m.s.)	&	male speaking (e.g. *bətaw ‘sister (m.s.)’)	\\
MSEA	&	Mainland South\ili{east} Asia	\\
N.	&	North (e.g. N. \ili{Wemale})	\\
(n.)	&	noun, nominal (e.g. sail (n.))	\\
NC	&	nasal-consonant cluster	\\
\tsc{neg}	&	negation, negative, no, not	\\
\tsc{neg1}	&	first part of bipartite negation	\\
\tsc{neg2}	&	second part of bipartite negation	\\
\tsc{nhum}	&	non-human 	\\
\ili{Nga'o}, Wat.	&	Watumite \ili{Nga'o}	\\
\ili{Nila} (W)	&	\ili{Nila} from Wotay village	\\
\ili{Nila} (K)	&	\ili{Nila} from Kokroman village	\\
\tsc{nmlz}	&	nominalizer	\\
non-An	&	non-Austronesian, unrelated Papuan language families	\\
\tsc{Oc}	&	Oceanic	\\
pers. comm.	&	personal communication	\\
P(CE)MP	&	PMP or PCEMP	\\
P(W)MP	&	PMP or PWMP	\\
PAMS	&	Proto-\ili{Ambon}, \citet{Stresemann1927} = PCM	\\
\tsc{PAn}	&	Proto-Austronesian	\\
\tsc{PAnD}	&	Proto-Austronesian = PMP, \citet{Dempwolff1938}	\\
\tsc{PAn-F}	&	Proto-Austronesian (Formosan)	\\
\tsc{PAnS}	&	Proto-Austronesian = PMP, \citet{Stresemann1927}	\\
PAS	&	Proto-\ili{Ambon}-Seram	\\
\it{PCEMP}	&	Proto-\ili{Central}-Eastern Malayo-Polynesian \it{(disputed node)}	\\
PCM	&	Proto-\ili{Central} Maluku	\\
\it{PCMP}	&	Proto-\ili{Central} Malayo-Polynesian \it{(disputed node)}	\\
%PECM	&	Proto-East \ili{Central} Maluku	\\
PEMP	&	Proto-Eastern Malayo-Polynesian	\\
\tsc{pfv}	&	perfective ‘already’	\\
PGWB	&	Proto-Greater West Bomberai \citep{UsherSchapper2022}	\\
\tsc{pi}	&	plural inclusive	\\
\tsc{pl}	&	plural 	\\
PMP	&	Proto-Malayo-Polynesian	\\
\tsc{Poc}	&	Proto-Oceanic	\\
\tsc{poss}	&	possessive	\\
PP	&	grandparent (parent's parent)	\\
\tsc{PPh}	&	Proto-Philippines	\\
PRM	&	Proto-Rote-\ili{Meto} \citep{Edwards2021}	\\
\tsc{prox}	&	proximal in time, space or reference; this, here, now	\\
%PSB	&	Proto-Sula-\ili{Buru}	\\
%PTAP	&	Proto-Timor-\ili{Alor}-\ili{Pantar} (Papuan)	\\
%PTNS	&	Proto-\ili{Teun}-\ili{Nila}-\ili{Serua}	\\
\it{PWMP}	&	Proto-Western Malayo-Polynesian \it{(disputed node)}	\\
\tsc{px}	&	plural exclusive	\\
\tsc{recip}	&	reciprocal	\\
S.	&	South (e.g. S. \ili{Wemale}, S. Halmahera)	\\
S\sub{A}	&	single argument of active intransitive clause (subject as Actor)	\\
SB	&	\tsc{Sula-Buru}	\\
SE.	&	South\ili{east} (e.g. SE Sulawesi)	\\
\tsc{sg}	&	singular 	\\
SHWNG	&	South Halmahera-West New Guinea	\\
SOV	&	Subject-Object-Verb	\\
sp.	&	species (singular; flora and fauna)	\\
spp.	&	species (plural; flora and fauna)	\\
\tsc{stat}	&	stative (BE relationship)	\\
\tsc{SSul}	&	South Sulawesi	\\
STB	&	\tsc{Seram-Tanimbar-Bomberai}	\\
S\sub{U}	&	single argument of non-active intransitive clause (subject as Undergoer)	\\
%Sul.	&	Sulawesi	\\
SVO	&	Subject-Verb-Object	\\
SW.	&	Southwest	\\
\tsc{Tal}	&	\ili{Taliabu}	\\
TAP	&	Timor-\ili{Alor}-\ili{Pantar} (Papuan language family)	\\
TB	&	\tsc{Timor-Babar}	\\
TNG	&	Trans-New Guinea (Papuan language family)	\\
TNS	&	\ili{Teun}-\ili{Nila}-\ili{Serua}	\\
\tsc{tr}	&	traces found (broader than \tsc{+fr})	\\
Uruang.	&	\ili{Uruangnirin}	\\
\tsc{u}	&	Undergoer, macrorole	\\
V	&	vowel	\\
(v.)	&	verb (e.g. `sail (v.)')	\\
vi.	&	intransitive verb	\\
VOC	&	Vereenigde Oost-Indische Compagnie (Dutch East Indies Trading Company)	\\
VSO	&	Verb-Subject-Object	\\
vt.	&	transitive verb	\\
VV	&	vowel sequence	\\
W.	&	west	\\
Watub.	&	\ili{Watubela}	\\
WB	&	wife’s brother	\\
\ili{Wetan} (T)	&	\ili{Wetan} from \ili{Teun} Island	\\
\ili{Wetan} (W)	&	\ili{Wetan} from \ili{Wetan} island	\\
Wey.	&	\ili{Weyewa}	\\
WMP	&	Western Malayo-Polynesian \it{(disputed node)}	\\
ySib	&	younger sibling	\\
ZH	&	sister’s husband	\\
	\lspbottomrule
\end{xltabular}}